\clearpage
\section{Conformance}\label{page:conformance}

\subsection{Conformance Claim}

A Bedrock implementation conformance claim identifies the architecture revision reported by CPUID and the CPUID
feature set for the logical processor to which the claim applies. For that revision and discovered feature set,
the implementation accepts and executes every architecturally valid instruction record according to its encoding,
state, event, memory, and instruction-specific rules. The same claim covers the architected discovery results,
save images, event frames, and reset state defined for that processor.

\subsection{Normative Material and Specificity}

Architecture prose, ordered steps, architecture tables, instruction encoding diagrams, and the Operation
and Detailed Semantics fields of instruction entries are normative unless a section explicitly identifies material
as a summary, example, or validation status.

\noindent
\begin{tabularx}{\linewidth}{@{}>{\raggedright\arraybackslash}p{1.55in}X@{}}
\toprule
\textbf{Rule set or provider} & \textbf{Covered material}\\
\midrule
Instruction definitions & Concrete instruction encodings, operand and effective-address grammar, stable form
identities, extension wiring, and document ordering.\\
\hyperref[page:instruction-execution-model]{Instruction Execution Model} common rules and instruction-entry Operation
and Detailed Semantics & Executable instruction behavior, architectural state transitions, fault and commit ordering,
repeat and event behavior, and single-logical-processor memory-action sequencing.\\
Formal memory model & Permitted cross-processor ordering and visibility.\\
ABI specifications & ELF, C ABI, and calling-convention contracts.\\
Compiler-interface specifications & Source-language and compiler-facing target-interface contracts.\\
External numerical floating-point providers & Numerical results and certificates only; input validation and the resulting
architectural trap or commit behavior follow the applicable
\hyperref[page:instruction-execution-model]{Instruction Execution Model} common rules and instruction entry's Operation
and Detailed Semantics.\\
\bottomrule
\end{tabularx}

Presented material remains normative where declared and follows the applicable
rule set. For executable behavior, readers apply the
\hyperref[page:instruction-execution-model]{Instruction Execution Model} common
rules together with the applicable instruction entry's Operation and Detailed
Semantics. The following specificity rules govern narrowing among applicable
rules:

\begin{enumerate}
\item Instruction bytes, field values, lengths, and decoding validity follow the concrete instruction form, its
field constraints, and the Instruction Encoding and Effective Addressing chapters.
\item Execution of a decoded instruction follows its Operation and Detailed Semantics. An instruction-specific rule
narrows the common execution, operand, flag, memory, or event rule that it names.
\item Common architectural behavior follows the chapter that defines that behavior. A structure-specific rule
narrows the defaults in Terminology and Compatibility.
\item Summary tables and examples provide navigation or explanatory context
and do not replace a normative rule.
\end{enumerate}

Two normative statements at the same specificity are required to agree. A discrepancy at that level is a
specification defect and is resolved by correcting the normative sources and their derived tables together.

\subsection{Implementation-Defined Disclosures}

An implementation-defined selection is stable for the implementation revision named by CPUID and is published
through the channel named by its defining rule. The publication identifies the selected value or behavior and the
processor revisions to which it applies.

\Needspace{3.0in}
\manualtablecaption{Implementation-Defined Disclosure Register}
\begin{manuallongtable}{@{}>{\raggedright\arraybackslash}p{1.35in}>{\raggedright\arraybackslash}p{2.20in}>{\raggedright\arraybackslash}p{2.10in}@{}}
\toprule
\rowcolor{ManualHeaderFill}
\textbf{Item} & \textbf{Defining rule} & \textbf{Publication}\\
\midrule
\endfirsthead
\multicolumn{3}{l}{\scriptsize\itshape Table \themanualtable\ (continued)}\\
\toprule
\rowcolor{ManualHeaderFill}
\textbf{Item} & \textbf{Defining rule} & \textbf{Publication}\\
\midrule
\endhead
@IMPLEMENTATION_DEFINED_ROWS@
\bottomrule
\end{manuallongtable}
